\section{Jejaring Kurikulum}

Jejaring kurikulum berikut menunjukkan keterkaitan antar mata kuliah pada Kurikulum 2025 Program Studi D4 Teknik Informatika: garis dari mata kuliah A ke B berarti B secara substansi membangun di atas materi A. Jalur bergaris hijau menandai mata kuliah umum/pengembangan diri, terpisah dari jalur teknis (biru); lihat legenda di bawah diagram untuk arti setiap gaya garis.

\definecolor{softgreen}{HTML}{D9F2DA}

\begin{figure}[H]
\centering
\resizebox{\textwidth}{!}{%
\begin{tikzpicture}[
  mk/.style={draw, rounded corners, fill=headblue, font=\fontsize{7}{8}\selectfont, align=center, text width=2.8cm, minimum height=0.8cm, inner sep=1.5pt},
  soft/.style={draw, rounded corners, fill=softgreen, font=\fontsize{7}{8}\selectfont, align=center, text width=2.8cm, minimum height=0.8cm, inner sep=1.5pt},
  gate/.style={draw, dashed, rounded corners, fill=headgray, font=\fontsize{7}{8}\selectfont, align=center, text width=2.8cm, minimum height=0.8cm, inner sep=1.5pt},
  smlabel/.style={font=\bfseries\small, anchor=east},
  ->, >={Stealth[length=1.5mm]}, thick, rounded corners=2pt,
]

% Semester row labels (mata kuliah dalam satu semester disusun berjenjang/ladder --
% posisi y berselang-seling -- bukan sejajar dalam satu garis, agar setiap kotak
% punya ruang vertikal sendiri dan lebih mudah ditata ulang tanpa saling tumpang tindih)
\node[smlabel] at (-2.0,0.000)       {Sem. 1};
\node[smlabel] at (-2.0,-4.970)    {Sem. 2};
\node[smlabel] at (-2.0,-11.340)   {Sem. 3};
\node[smlabel] at (-2.0,-17.710)   {Sem. 4};
\node[smlabel] at (-2.0,-24.510)   {Sem. 5};
\node[smlabel] at (-2.0,-29.080)   {Sem. 6};
\node[smlabel] at (-2.0,-44.960)   {Sem. 7};
\node[smlabel] at (-2.0,-49.170)   {Sem. 8};

% ===================== SEMESTER 1 =====================
\node[mk]   (n251002) at (0.000,0.000)   {RTI251002\\Konsep TI};
\node[mk]   (n251004) at (3.300,-1.150)  {RTI251004\\Matematika Dasar};
\node[mk]   (n251006) at (6.600,0.000)   {RTI251006\\Dasar Pemrograman};
\node[mk]   (n251007) at (9.900,-1.150)  {RTI251007\\Prak. Dasar Pemrograman};
\node[mk]   (n251009) at (13.200,0.000)  {RTI251009\\Fisika};
\node[soft] (n251005) at (16.500,-1.150) {RTI251005\\Bahasa Inggris 1};
\node[soft] (n251001) at (19.800,0.000)  {RTI251001\\Pancasila};
\node[soft] (n251008) at (23.100,-1.150) {RTI251008\\K3};
\node[soft] (n251003) at (26.400,0.000)  {RTI251003\\Critical Thinking};

% ===================== SEMESTER 2 =====================
\node[mk]   (n252008) at (0.000,-4.970)   {RTI252008\\Algoritma \& Struktur Data};
\node[mk]   (n252009) at (3.300,-6.120)   {RTI252009\\Prak. Alg. \& Struktur Data};
\node[mk]   (n252002) at (6.600,-4.970)   {RTI252002\\Aljabar Linier};
\node[mk]   (n252003) at (9.900,-6.120)   {RTI252003\\Desain Antarmuka};
\node[mk]   (n252004) at (13.200,-4.970)  {RTI252004\\Sistem Operasi};
\node[mk]   (n252005) at (16.500,-6.120)  {RTI252005\\Rekayasa Perangkat Lunak};
\node[mk]   (n252006) at (19.800,-4.970)  {RTI252006\\Basis Data};
\node[mk]   (n252007) at (23.100,-6.120)  {RTI252007\\Prak. Basis Data};
\node[soft] (n252001) at (26.400,-4.970)  {RTI252001\\Agama};

% ===================== SEMESTER 3 =====================
\node[mk]   (n253006) at (0.000,-11.340)  {RTI253006\\Metode Numerik};
\node[mk]   (n253007) at (3.300,-12.490)  {RTI253007\\Pemrog. Berbasis Objek};
\node[mk]   (n253008) at (6.600,-11.340)  {RTI253008\\Prak. Pemrog. Berbasis Objek};
\node[mk]   (n253001) at (9.900,-12.490)  {RTI253001\\Manajemen Proyek};
\node[mk]   (n253004) at (13.200,-11.340) {RTI253004\\Desain \& Pemrog. Web};
\node[mk]   (n253002) at (16.500,-12.490) {RTI253002\\Sistem Informasi Manajemen};
\node[mk]   (n253005) at (19.800,-11.340) {RTI253005\\Basis Data Lanjut};
\node[soft] (n253003) at (23.100,-12.490) {RTI253003\\Kewarganegaraan};
\node[soft] (n253009) at (26.400,-11.340) {RTI253009\\Bahasa Inggris 2};

% ===================== SEMESTER 4 =====================
\node[mk]   (n254001) at (0.000,-17.710)  {RTI254001\\Sistem Pendukung Keputusan};
\node[mk]   (n254004) at (3.300,-18.860)  {RTI254004\\Kecerdasan Artifisial};
\node[mk]   (n254002) at (6.600,-17.710)  {RTI254002\\Analisis \& Desain Berorientasi Objek};
\node[mk]   (n254009) at (9.900,-18.860)  {RTI254009\\Proyek Sistem Informasi};
\node[mk]   (n254005) at (13.200,-17.710) {RTI254005\\Jaringan Komputer};
\node[mk]   (n254006) at (16.500,-18.860) {RTI254006\\Prak. Jaringan Komputer};
\node[mk]   (n254007) at (19.800,-17.710) {RTI254007\\Pemrograman Web Lanjut};
\node[mk]   (n254008) at (23.100,-18.860) {RTI254008\\Statistik Komputasi};
\node[soft] (n254003) at (26.400,-17.710) {RTI254003\\Bahasa Indonesia};

% ===================== SEMESTER 5 =====================
\node[mk]   (n255003) at (0.000,-24.510)  {RTI255003\\Pemrograman Mobile};
\node[mk]   (n255004) at (3.300,-25.660)  {RTI255004\\Pembelajaran Mesin};
\node[mk]   (n255007) at (6.600,-24.510)  {RTI255007\\Pengolahan Citra \& Visi Komputer};
\node[mk]   (n255002) at (9.900,-25.660)  {RTI255002\\Metodologi Penelitian};
\node[mk]   (n255005) at (13.200,-24.510) {RTI255005\\Business Intelligence};
\node[mk]   (n255006) at (16.500,-25.660) {RTI255006\\Penjaminan Mutu PL};
\node[mk]   (n255008) at (19.800,-24.510) {RTI255008\\Adm. \& Keamanan Jaringan};
\node[soft] (n255001) at (23.100,-25.660) {RTI255001\\Kewirausahaan Berbasis Teknologi};

% ===================== SEMESTER 6 =====================
\node[mk]   (n256003) at (3.300,-29.080)  {RTI256003\\Pemrog. Berbasis Framework};
\node[soft] (n256001) at (6.600,-30.230)  {RTI256001\\Komunikasi \& Etika Profesi};
\node[soft] (n256002) at (14.850,-29.080) {RTI256002\\Pengembangan Karir};

\node[font=\bfseries\small] at (3.3,-31.480)   {Jalur Reguler};
\node[font=\bfseries\small] at (14.85,-31.480) {Jalur Magang};
\draw[densely dotted, thick, -] (9.9,-28.320) -- (9.9,-42.560);

\node[mk] (n256104) at (0.000,-35.440)  {RTI256104\\Proyek Tek. Terintegrasi};
\node[mk] (n256105) at (3.300,-36.590)  {RTI256105\\Internet of Things};
\node[mk] (n256106) at (6.600,-35.440)  {RTI256106\\Big Data};
\node[mk] (n256107) at (0.000,-40.650)  {RTI256107\\Cloud Computing};
\node[mk] (n256108) at (3.300,-41.800)  {RTI256108\\Komputasi Hijau};

\node[mk] (n256204) at (13.200,-36.590) {RTI256204\\Proyek Inovasi};
\node[mk] (n256205) at (16.500,-35.440) {RTI256205\\Workshop Tek. Terapan};
\node[mk] (n256206) at (13.200,-40.650) {RTI256206\\Rekayasa Sistem};
\node[mk] (n256207) at (16.500,-41.800) {RTI256207\\Teknologi Terapan};

% ===================== SEMESTER 7 =====================
\node[gate] (gate120) at (3.300,-44.960) {Lulus Semester 1--6\\($\geq$120 SKS)};
\node[mk]   (n257001) at (9.900,-46.110) {RTI257001\\Magang};

% ===================== SEMESTER 8 =====================
\node[gate] (gate130) at (3.300,-49.170) {$\geq$130 SKS};
\node[mk]   (n258001) at (9.900,-50.320) {RTI258001\\Skripsi};
\node[soft] (n258002) at (16.500,-49.170) {RTI258002\\Bahasa Inggris Persiapan Kerja};

% ===================== EDGES: JALUR TEKNIS, PERLU VALIDASI, DAN UMUM =====================
% Semua jalur digambar siku-siku (bukan garis lurus/diagonal). Jalur dalam satu
% semester memakai lekukan-S melalui celah antar level ladder; jalur antar semester
% keluar ke celah di bawah baris asal, lalu menuju jalur (lane) tegak yang bebas dari
% kotak manapun -- diprioritaskan jalur di sela-sela kolom (paling dekat/pendek),
% baru jalur di pinggir kiri/kanan grid bila semua sela kolom terpakai -- sebelum
% masuk kembali dari atas baris tujuan. Posisi setiap jalur (shelf keluar/masuk dan
% jalur tegak) dipilih agar tidak memotong kotak mata kuliah lain DAN tidak berimpit
% dengan jalur lain pada jarak cetak sebenarnya (diverifikasi terhadap skala cetak
% buku ini, bukan hanya satuan tikz).
\draw (n251006.east) -- (8.250,0.000) -- (8.250,-1.150) -- (n251007.west);
\draw (n251006.south) -- (6.600,-0.660) -- (4.950,-0.660) -- (4.950,-4.310) -- (0.000,-4.310) -- (n252008.north);
\draw (n251004.south) -- (3.300,-1.810) -- (1.650,-1.810) -- (1.650,-10.680) -- (0.000,-10.680) -- (n253006.north);
\draw (n251006.south) -- (6.600,-0.660) -- (8.250,-0.660) -- (8.250,-10.730) -- (3.300,-10.730) -- (n253007.north);
\draw (n251006.south) -- (6.600,-1.710) -- (11.550,-1.710) -- (11.550,-10.680) -- (6.600,-10.680) -- (n253008.north);
\draw (n253007.east) -- (4.950,-12.490) -- (4.950,-11.340) -- (n253008.west);
\draw (n251006.south) -- (6.600,-1.710) -- (14.850,-1.710) -- (14.850,-10.680) -- (13.200,-10.680) -- (n253004.north);
\draw (n252006.south) -- (19.800,-5.630) -- (18.150,-5.630) -- (18.150,-11.830) -- (16.500,-11.830) -- (n253002.north);
\draw (n252006.south) -- (19.800,-5.630) -- (19.800,-10.680) -- (n253005.north);
\draw (n252006.south) -- (19.800,-5.630) -- (21.450,-5.630) -- (21.450,-10.680) -- (13.200,-10.680) -- (n253004.north);
\draw (n252008.south) -- (0.000,-6.680) -- (24.750,-6.680) -- (24.750,-17.100) -- (3.300,-17.100) -- (n254004.north);
\draw (n252005.south) -- (16.500,-7.180) -- (32.000,-7.180) -- (32.000,-16.600) -- (6.600,-16.600) -- (n254002.north);
\draw (n252005.south) -- (16.500,-7.180) -- (32.500,-7.180) -- (32.500,-16.600) -- (9.900,-16.600) -- (n254009.north);
\draw (n252005.south) -- (16.500,-7.180) -- (33.000,-7.180) -- (33.000,-23.900) -- (16.500,-23.900) -- (n255006.north);
\draw (n253004.south) -- (13.200,-13.100) -- (4.950,-13.100) -- (4.950,-16.100) -- (19.800,-16.100) -- (n254007.north);
\draw (n253001.south) -- (9.900,-13.150) -- (9.900,-18.200) -- (n254009.north);
\draw[dashed] (n254005.east) -- (14.850,-17.710) -- (14.850,-18.860) -- (n254006.west);
\draw (n252004.south) -- (13.200,-7.680) -- (33.500,-7.680) -- (33.500,-23.400) -- (19.800,-23.400) -- (n255008.north);
\draw (n252003.south) -- (9.900,-7.180) -- (-4.000,-7.180) -- (-4.000,-23.850) -- (0.000,-23.850) -- (n255003.north);
\draw (n252002.south) -- (6.600,-7.680) -- (-4.500,-7.680) -- (-4.500,-23.350) -- (3.300,-23.350) -- (n255004.north);
\draw (n252006.south) -- (19.800,-8.180) -- (34.000,-8.180) -- (34.000,-22.900) -- (13.200,-22.900) -- (n255005.north);
\draw (n254005.south) -- (13.200,-18.370) -- (11.550,-18.370) -- (11.550,-23.400) -- (19.800,-23.400) -- (n255008.north);
\draw (n254008.south) -- (23.100,-19.520) -- (21.450,-19.520) -- (21.450,-22.400) -- (9.900,-22.400) -- (n255002.north);
\draw (n254008.south) -- (23.100,-19.520) -- (26.400,-19.520) -- (26.400,-22.400) -- (3.300,-22.400) -- (n255004.north);
\draw (n255004.east) -- (4.950,-25.660) -- (4.950,-24.510) -- (n255007.west);
\draw (n254007.south) -- (19.800,-18.370) -- (18.150,-18.370) -- (18.150,-28.420) -- (3.300,-28.420) -- (n256003.north);
\draw (n254005.south) -- (13.200,-19.470) -- (8.250,-19.470) -- (8.250,-34.830) -- (3.300,-34.830) -- (n256105.north);
\draw (n254005.south) -- (13.200,-19.470) -- (1.650,-19.470) -- (1.650,-39.990) -- (0.000,-39.990) -- (n256107.north);
\draw (n255008.south) -- (19.800,-25.170) -- (19.800,-39.990) -- (0.000,-39.990) -- (n256107.north);
\draw (n254009.south) -- (9.900,-19.970) -- (-5.000,-19.970) -- (-5.000,-34.780) -- (0.000,-34.780) -- (n256104.north);
\draw (n256003.south) -- (3.300,-29.740) -- (3.300,-34.780) -- (0.000,-34.780) -- (n256104.north);
\draw (n254009.south) -- (9.900,-19.970) -- (-5.500,-19.970) -- (-5.500,-34.280) -- (13.200,-34.280) -- (n256204.north);
\draw (n255001.south) -- (23.100,-26.320) -- (23.100,-34.830) -- (13.200,-34.830) -- (n256204.north);
\draw (n254009.south) -- (9.900,-19.970) -- (-6.000,-19.970) -- (-6.000,-34.280) -- (16.500,-34.280) -- (n256205.north);
\draw (n255006.south) -- (16.500,-26.320) -- (16.500,-34.780) -- (n256205.north);
\draw (n254009.south) -- (9.900,-19.970) -- (-6.500,-19.970) -- (-6.500,-39.490) -- (13.200,-39.490) -- (n256206.north);
\draw (n254002.south) -- (6.600,-20.470) -- (-7.000,-20.470) -- (-7.000,-39.490) -- (13.200,-39.490) -- (n256206.north);
\draw (n254009.south) -- (9.900,-19.970) -- (-7.500,-19.970) -- (-7.500,-38.990) -- (16.500,-38.990) -- (n256207.north);
\draw (n253001.south) -- (9.900,-13.600) -- (-8.000,-13.600) -- (-8.000,-38.990) -- (16.500,-38.990) -- (n256207.north);
\draw (gate120.east) -- (6.600,-44.960) -- (6.600,-46.110) -- (n257001.west);
\draw (gate130.east) -- (6.600,-49.170) -- (6.600,-50.320) -- (n258001.west);
\draw (n255002.south) -- (9.900,-26.320) -- (4.950,-26.320) -- (4.950,-49.660) -- (9.900,-49.660) -- (n258001.north);
\draw[dotted] (n255005.south) -- (13.200,-25.170) -- (13.200,-33.780) -- (6.600,-33.780) -- (n256106.north);
\draw[dotted] (n256107.east) -- (1.650,-40.650) -- (1.650,-41.800) -- (n256108.west);
\draw[dotted] (n254004.south) -- (3.300,-19.520) -- (3.300,-25.000) -- (n255004.north);
\draw[dotted] (n253007.south) -- (3.300,-14.100) -- (-8.500,-14.100) -- (-8.500,-23.850) -- (0.000,-23.850) -- (n255003.north);
\draw[green!40!black] (n251005.south) -- (16.500,-1.810) -- (34.500,-1.810) -- (34.500,-10.680) -- (26.400,-10.680) -- (n253009.north);
\draw[green!40!black] (n253009.south) -- (26.400,-12.000) -- (35.000,-12.000) -- (35.000,-48.510) -- (16.500,-48.510) -- (n258002.north);
\draw[green!40!black] (n251001.south) -- (19.800,-2.310) -- (35.500,-2.310) -- (35.500,-10.180) -- (23.100,-10.180) -- (n253003.north);
\draw[green!40!black] (n256001.south) -- (6.600,-30.890) -- (9.900,-30.890) -- (9.900,-45.450) -- (n257001.north);
\draw[green!40!black] (n256002.south) -- (14.850,-29.740) -- (14.850,-45.450) -- (9.900,-45.450) -- (n257001.north);

\end{tikzpicture}%
}

\vspace{2mm}
\resizebox{0.85\textwidth}{!}{%
\begin{tikzpicture}[font=\scriptsize, ->, >={Stealth[length=1.5mm]}, thick]
\draw (0,0) -- (0.8,0); \node[anchor=west] at (0.9,0) {Keterkaitan konsep terkonfirmasi};
\draw[dotted] (6.7,0) -- (7.5,0); \node[anchor=west] at (7.6,0) {Perlu validasi (kesamaan topik)};
\draw[dashed] (13.3,0) -- (14.1,0); \node[anchor=west] at (14.2,0) {Korekuisit (satu semester)};
\draw[green!40!black] (18.5,0) -- (19.3,0); \node[anchor=west] at (19.4,0) {Jalur umum/pengembangan diri};
\end{tikzpicture}%
}
\caption{Jejaring Kurikulum Program Studi D4 Teknik Informatika (Kurikulum 2025)}
\end{figure}

\reviewfrompdf{Diagram jejaring kurikulum di atas adalah visualisasi baru yang disusun dari data Deskripsi Mata Kuliah, Bahan Kajian, dan Matakuliah Syarat pada RPS Kurikulum 2025, direpresentasikan sebagai peta keterkaitan konsep (bukan prasyarat pendaftaran) mengingat mahasiswa menempuh satu paket mata kuliah tetap per semester (kecuali percabangan Jalur Reguler/Jalur Magang pada Semester 6); dokumen kurikulum 2020 tidak memiliki data digital yang setara untuk Peta Kurikulum dan Pohon Kurikulum (lihat catatan pada bagian berikut). Sejumlah keterkaitan (ditandai garis bertitik) disusun berdasarkan kesamaan topik/bahan kajian, bukan pernyataan eksplisit pada RPS, dan perlu divalidasi oleh tim kurikulum.}

\clearpage
\restoregeometry

